{\scriptsize
\setlength{\tabcolsep}{3pt}
\begin{longtable}{>{\raggedright\arraybackslash}p{2.55cm}>{\raggedright\arraybackslash}p{2.55cm}>{\raggedright\arraybackslash}p{2.55cm}>{\raggedright\arraybackslash}p{2.55cm}>{\raggedright\arraybackslash}p{2.55cm}>{\raggedright\arraybackslash}p{2.55cm}}
\caption{Nombre de composantes actives du VB selon $\alpha_0$ et le type de covariance ($K_{\max}=6$)}\label{tab:S8}\\
\toprule
\textbf{Covariance} & \textbf{$\alpha_0$} & \textbf{$\bar K$ actif (1 \%)} & \textbf{Écart-type} & \textbf{ARI} & \textbf{Non convergés} \\
\midrule
\endfirsthead
\toprule
\textbf{Covariance} & \textbf{$\alpha_0$} & \textbf{$\bar K$ actif (1 \%)} & \textbf{Écart-type} & \textbf{ARI} & \textbf{Non convergés} \\
\midrule
\endhead
\bottomrule
\endfoot
diag & 1,00 & 3,33 & 0,76 & 0,942 & 0\,\% \\
diag & 1,50 & 3,87 & 0,94 & 0,916 & 0\,\% \\
diag & 2,00 & 4,30 & 1,02 & 0,901 & 0\,\% \\
diag & 2,50 & 4,90 & 1,03 & 0,890 & 0\,\% \\
diag & 3,00 & 5,13 & 0,90 & 0,897 & 0\,\% \\
diag & 5,00 & 6,00 & 0,00 & 0,826 & 0\,\% \\
diag & 10,00 & 6,00 & 0,00 & 0,732 & 0\,\% \\
full & 1,00 & 3,00 & 0,00 & 0,961 & 0\,\% \\
full & 1,50 & 3,03 & 0,18 & 0,961 & 0\,\% \\
full & 2,00 & 3,03 & 0,18 & 0,961 & 0\,\% \\
full & 2,50 & 3,13 & 0,35 & 0,956 & 0\,\% \\
full & 3,00 & 3,33 & 0,48 & 0,951 & 0\,\% \\
full & 5,00 & 6,00 & 0,00 & 0,929 & 0\,\% \\
full & 10,00 & 6,00 & 0,00 & 0,724 & 0\,\% \\
spherical & 1,00 & 3,67 & 0,92 & 0,935 & 0\,\% \\
spherical & 1,50 & 4,40 & 1,07 & 0,909 & 0\,\% \\
spherical & 2,00 & 4,77 & 1,14 & 0,881 & 0\,\% \\
spherical & 2,50 & 5,67 & 0,66 & 0,854 & 0\,\% \\
spherical & 3,00 & 5,97 & 0,18 & 0,837 & 0\,\% \\
spherical & 5,00 & 6,00 & 0,00 & 0,796 & 0\,\% \\
spherical & 10,00 & 6,00 & 0,00 & 0,676 & 0\,\% \\
\end{longtable}
\par\noindent{\footnotesize Seuils théoriques de Rousseau et Mengersen (2011) en $D=2$ : $d/2=2{,}5$ (pleine), $2$ (diagonale), $1{,}5$ (sphérique) ; hypothèse à tester, la théorie portant sur le posterior exact. Covariances tied/diag/spherical : critère approché.}
}
